\chapter{Difetti e criticità del sistema}\label{chap:defects}


\noindent Questo capitolo raccoglie i difetti individuati nel corso dell'analisi. Essi non sono
un residuo accessorio del lavoro: la loro individuazione è stata resa necessaria dalla necessità
di stabilire se i risultati negativi fossero da attribuire alle strategie sviluppate o al
materiale di partenza, e la distinzione si è rivelata decisiva in più di un caso. Il capitolo
separa i difetti del sistema studiato da quelli del codice sviluppato in questo lavoro --- i primi
limiti del materiale, i secondi questione di attendibilità delle misure --- ne quantifica
l'impatto, ne indica il rimedio e ne trae le conseguenze generali per chi impiegasse il sistema su
altri dati.

% =======================================================================================
\section{Classificazione dei difetti}\label{sec:defects-taxonomy}

I difetti individuati appartengono a quattro categorie, che è opportuno tenere distinte perché
richiedono interventi di natura diversa. La distinzione che più conta, ai fini della lettura del
capitolo, è però un'altra: quella fra i difetti del \emph{sistema studiato} --- quattro, tutti nel
archivio --- e gli errori del \emph{codice sviluppato in questo lavoro} --- due,
entrambi individuati e corretti durante l'analisi. I primi limitano ciò che si può concludere sul
sistema; i secondi riguardano l'attendibilità delle misure e sono esposti perché il loro
trattamento fa parte del metodo.

\begin{table}[htbp]
\centering
\caption{Difetti individuati, per categoria e provenienza.}
\label{tab:defects}
\footnotesize
\begin{tabular}{@{}clll@{}}
\toprule
\# & difetto & categoria & provenienza \\
\midrule
D1 & ramo irraggiungibile nel calcolo della similarità & algoritmo & materiale d'archivio \\
D2 & disallineamento di schema nei dataset Glass & dati & materiale d'archivio \\
D3 & indici di ground truth non ricostruibili (Physician) & dati & materiale d'archivio \\
D4 & eseguibile di build contenente istruzioni di debug & artefatto & materiale d'archivio \\
D5 & computo errato del budget nella politica sui veti & codice sviluppato & \textbf{questo lavoro} \\
D6 & indice di rischio anti-correlato con la rilevanza & codice sviluppato & \textbf{questo lavoro} \\
\bottomrule
\end{tabular}
\end{table}

% =======================================================================================
\begin{table}[htbp]
\centering
\caption{Impatto dei difetti, come sono stati rilevati e stato attuale. Le classi sono quelle
della \autoref{tab:defects}.}
\label{tab:defects-impact}
\footnotesize
\begin{tabular}{@{}lp{4.1cm}p{3.1cm}p{2.1cm}@{}}
\toprule
difetto & impatto quantificato & come è stato rilevato & stato \\
\midrule
D1 & la conversione \texttt{B}$\to$\texttt{D} è peggiore in \num{25} configurazioni su \num{25},
con una perdita di \num{21.7} punti di accuratezza & confronto fra i due rami sul bytecode, poi
misura & non corretto: il vincolo è di non modificare l'eseguibile \\
D2 & la famiglia Glass non è utilizzabile come distribuita: lo schema dei dataset è posizionale,
quello dei file di regole usa i nomi & prima esecuzione fallita, poi ispezione delle intestazioni & corretto esternamente, con rinomina documentata \\
D3 & due configurazioni non eseguibili per righe malformate; \num{115} celle non localizzabili su
\num{404}; celle valutabili da \num{329} a \num{40} & esecuzione fallita per eccezione di
indice fuori intervallo; conteggio delle righe con campi mancanti & non riparato: costo di preprocessing non giustificato \\
D4 & una build alternativa con due istruzioni di diagnostica, con comportamento identico
sull'output & confronto fra i due eseguibili sulla stessa configurazione & documentato: esclude che l'archivio contenga una versione emendata di D1 \\
D5 & riduzione allo \SI{0.1}{\percent} invece del \SI{17.1}{\percent} su Bridges a soglia 9;
\SI{43.7}{\percent} delle rimozioni dichiarate esatte senza copritore sopravvissuto & confronto fra
riduzione dichiarata ed effettiva su una pipeline composita & corretto e verificato su \num{8} configurazioni su \num{8} \\
D6 & indice di rischio con $\rho = -0.237$ su \num{8961} regole; dopo la correzione, $2.5$ volte
più regole inutili rimosse e $2.6$ volte meno celle servibili distrutte & validazione dell'indice
contro il numero di celle che ciascuna regola può servire & corretto, senza guadagno sull'esito \\
\bottomrule
\end{tabular}
\end{table}

\section{Difetti del sistema studiato}\label{sec:defects-system}

I quattro difetti che seguono appartengono al materiale d'archivio del sistema studiato. Essi hanno
una caratteristica comune: incidono sulla \emph{possibilità} di interpretare i risultati, non
soltanto sulla loro qualità, e in due casi hanno reso inutilizzabile una parte della griglia
sperimentale. Ciascuno è presentato con il sintomo, l'evidenza che lo stabilisce, l'impatto sui
risultati e il rimedio adottato.

\subsection{D1 --- Ramo irraggiungibile nel calcolo della similarità}\label{sec:d1}

Il sistema calcola il contributo di similarità di un attributo con una procedura che distingue
fra attributi numerici e attributi categorici. Il disassemblaggio mostra che la diramazione verso
il caso categorico è \emph{irraggiungibile}: la condizione di selezione confronta riferimenti a
oggetti anziché il contenuto, e non può quindi risultare vera. In conseguenza di ciò, un attributo
categorico viene elaborato dalla procedura numerica, che tenta una conversione a numero, fallisce,
e ricade su una distanza di stringa applicata a etichette che possono essere lunghe e prive di
relazione.

\begin{observation}
Il difetto è confermato dall'analisi statica e non è stato corretto: il vincolo di non modificare
l'eseguibile è stato mantenuto per tutta la durata della valutazione esterna. La sua esistenza
non è un'ipotesi ma un fatto verificabile nel bytecode.
\end{observation}

\paragraph{Impatto sull'analisi.} Il difetto ha conseguenze sulla tipizzazione delle colonne
impiegata nelle campagne. Si è verificato sperimentalmente che la conversione di una colonna da
categoria a numerica, lungi dal costituire un miglioramento gratuito, peggiora il risultato:
su \num{25} configurazioni, la conversione produce una perdita media di \SI{21.7}{} punti di
accuratezza, con un peggioramento in \num{25} configurazioni su \num{25}. Il risultato ha
richiesto di rivedere una conclusione formulata in una fase precedente del lavoro, nella quale la
conversione era stata presentata come un guadagno.

% =======================================================================================
\subsection{D2 --- Disallineamento di schema nei dataset Glass}\label{sec:d2}

I dataset della famiglia Glass presentano nomi di colonna in forma posizionale
(\texttt{A}, \texttt{B}, \dots), mentre i corrispondenti file di regole impiegano nomi
descrittivi. Per le altre famiglie i due insiemi di nomi coincidono; per Glass no.

Il disallineamento non è innocuo. Il sistema risolve le colonne \emph{per posizione}, non per
nome: quando gli schemi non coincidono, la corrispondenza fra le colonne del dataset e quelle
delle regole è apparentemente corretta ma sostanzialmente errata, e il sistema scrive le
imputazioni nella colonna sbagliata. La condizione è di corruzione silenziosa: nessun errore
viene segnalato.

\paragraph{Intervento.} La correzione adottata è esterna e consiste nel rinominare
posizionalmente le colonne del dataset perché coincidano con quelle dei file di regole, e nel
riscrivere di conseguenza il materiale di ground truth. L'ordine delle colonne non è modificato,
sicché gli indici posizionali impiegati dal sistema restano validi.

\paragraph{Impatto sull'analisi.} Prima della correzione la famiglia Glass risultava
inutilizzabile, con copertura nulla. Dopo la correzione la copertura si attesta intorno al
\SI{30}{\percent}. Va osservato che il difetto ha un effetto anche sull'interpretazione delle
statistiche di training: i file di risultato prodotti prima della correzione registrano valori
non imputati per tutte le celle, e non sono quindi impiegabili come riferimento. La circostanza ha
richiesto di raccogliere statistiche ex novo su esecuzioni condotte con lo schema corretto.

% =======================================================================================
\subsection{D3 --- Righe malformate e indici non ricostruibili nei dataset Physician}
\label{sec:d3}

La famiglia Physician comprende quattro versioni ``remodulate'' di dimensioni contenute, pensate
per la sperimentazione. Il loro materiale presenta due difetti distinti, che una verifica
successiva ha permesso di separare.

\paragraph{Righe con un numero di campi inferiore allo schema.} I dataset delle versioni 0.05 e
0.1 contengono righe malformate: \num{7}, \num{11} o \num{1} campi invece di \num{13}. Il file
della versione 0.05 termina a metà di un record, e nella 0.1 una riga mostra due record
concatenati, con un valore testuale che risulta tagliato e ricomposto (``\texttt{PENNSYLVANIA
COLLEGSICAL THERAPY}''). È questa la causa diretta della non eseguibilità: il controllo di chiave
che il sistema esegue all'avvio accede ai campi di una tupla per indice, e su una riga corta esce
con \texttt{IndexOutOfBoundsException}. La diagnosi è confermata dal materiale d'archivio: per
entrambe le versioni il log, il file delle celle popolate e il file dei risultati spediti con
l'archivio sono \emph{vuoti}, mentre per le versioni 0.005 e 0.01 sono popolati.

\paragraph{Indici di ground truth eccedenti il dataset.} Il materiale di ground truth --- in
formato \texttt{riga;attributo;valore} --- contiene indici di riga che eccedono il numero di righe
del dataset corrispondente, ed è ben formato: il difetto riguarda il \emph{riferimento}, non la
sintassi.

\begin{table}[htbp]
\centering
\caption{Difetto D3. Righe malformate e voci di ground truth con indice fuori intervallo, per
versione remodulata. Le colonne ``valutabili'' indicano le celle che restano dopo l'esclusione
degli indici fuori intervallo.}
\label{tab:physician-gt}
\footnotesize
\begin{tabular}{lrrrrr}
\toprule
versione & righe & righe malformate & voci ground truth & fuori intervallo & valutabili \\
\midrule
0.005 & 103  & 0 & 13  & 0   & 13 \\
0.01  & 207  & 0 & 27  & 0   & 27 \\
0.05  & 914  & \textbf{1} & 135 & \textbf{14}  & \textbf{121} \\
0.1   & 1374 & \textbf{2} & 269 & \textbf{101} & \textbf{168} \\
\bottomrule
\end{tabular}
\end{table}

Gli indici fuori intervallo provengono dal dataset completo, prima della riduzione, e non sono
stati rimappati; le righe malformate indicano che la riduzione è avvenuta per \emph{troncamento}
del file, circostanza che spiega anche perché gli indici non corrispondano più. Escludendo le voci
non localizzabili, le celle valutabili passano da \num{40} a \num{329}: la correzione non richiede
di modificare il dataset, ma soltanto di dichiarare quali celle si escludono.

\paragraph{La ricostruzione della corrispondenza non è praticabile.} Si è verificato se gli indici
potessero essere ricostruiti confrontando le righe delle versioni ridotte con quelle del dataset
completo. La verifica ha dato esito negativo: su \num{200} righe esaminate, \num{29} non trovano
corrispondenza e \num{73} risultano ambigue. Le versioni remodulate non sono un sottocampionamento
del dataset completo ma una sua perturbazione, e l'informazione necessaria alla rimappatura non è
contenuta nei dati disponibili.

\paragraph{Riparazione possibile, non intrapresa.} Le righe malformate possono essere sostituite
da righe di soli valori mancanti, conservando la numerazione su cui il ground truth è indicizzato
e producendo righe che il sistema ignora in entrambi i ruoli; la procedura è implementata e
documentata (\texttt{pruning/repair\_physician\_datasets.py}). La riparazione non è stata
applicata alla campagna, perché il costo di preprocessing del sistema su queste configurazioni ---
il controllo di chiave è quadratico nel numero di righe, e vale decine di minuti per esecuzione
--- non è giustificato dal numero di celle che se ne ricaverebbero, e perché il contributo
scientifico delle due versioni sarebbe comunque marginale rispetto alle versioni già eseguibili.

\begin{observation}
La famiglia Physician è utilizzabile come quarta famiglia per la verifica di equivalenza, e lo è
stata (\autoref{sec:physician}); per la verifica di qualità fornisce \num{40} celle con copertura
non nulla, un campione insufficiente a trarre conclusioni. La limitazione è del materiale, non
dell'impostazione sperimentale, e va dichiarata come tale.
\end{observation}

% =======================================================================================
\subsection{D4 --- Eseguibile di build con istruzioni di debug}\label{sec:d4}

L'archivio contiene un secondo eseguibile, di dimensioni ridotte, la cui struttura interna è
danneggiata (la directory centrale dell'archivio è assente). L'estrazione delle voci è comunque
possibile, e il confronto con l'eseguibile principale mostra una sola differenza: la presenza di
due istruzioni di stampa a scopo di diagnostica. Il comportamento è per il resto identico, come
verificato eseguendo entrambi gli eseguibili sulla medesima configurazione e confrontando gli
output.

Non si tratta quindi di una versione con correzioni, ma di una build di sviluppo. La circostanza
è rilevante in senso negativo: esclude che l'archivio contenga una versione emendata del difetto
D1.

% =======================================================================================
\section{Errori rilevati e corretti nel codice sviluppato}\label{sec:defects-ours}

I due difetti che seguono non appartengono al sistema studiato ma agli strumenti sviluppati in
questo lavoro. La loro esposizione è deliberata e costituisce parte del contributo: entrambi sono
stati individuati da verifiche interne --- il confronto fra riduzione dichiarata e riduzione
effettiva, la validazione di un indice di rischio mai messo alla prova --- e in un caso la
correzione ha cambiato il valore di una misura riportata in una fase precedente dell'analisi. Un
lavoro che misuri l'effetto di una trasformazione su un sistema opaco deve poter esibire il
controllo di qualità dei propri strumenti, perché è da lì che provengono gli errori più difficili
da vedere.

\subsection{D5 --- Computo errato del budget nella politica sui veti}\label{sec:d5}

Questo difetto appartiene al codice sviluppato in questo lavoro, ed è esposto perché la sua
individuazione e correzione costituiscono parte del contributo.

\paragraph{Il difetto.} La politica \str{budget:F} limita la frazione di regole portatrici di veto
che possono essere rimosse. Il computo della frazione includeva però anche le rimozioni
\emph{provate prive di costo} dalla condizione \Vtwo{}: regole il cui potere di veto è implicato
da una regola sopravvissuta. Tali rimozioni non sacrificano alcun veto, e addebitarle costituiva
una doppia contabilizzazione.

\paragraph{Secondo difetto, nella stessa procedura.} La condizione \Vtwo{} veniva valutata
sull'insieme di regole \emph{in ingresso} anziché su quello \emph{sopravvissuto}. Se una strategia
successiva rimuove la regola coprente, la garanzia decade. Su una pipeline composita si è
misurato che il \SI{43.7}{\percent} delle rimozioni presentate come esatte non aveva più alcun
copritore sopravvissuto.

\paragraph{Un terzo effetto, nella politica di tutela.} La politica \str{protect}, che deve
ripristinare le regole portatrici di veto non altrimenti garantite, ripristinava anche le regole
la cui rimozione era dimostrabilmente priva di costo, annullando di fatto il risultato del
teorema. Sulla configurazione esaminata la riduzione scendeva allo \SI{0.1}{\percent} invece del
\SI{17.1}{\percent} che la condizione autorizza.

\paragraph{Correzione e verifica.} La condizione \Vtwo{} è stata riformulata in modo da essere
valutata sull'insieme corrente delle regole sopravvissute, e il computo del budget è stato
corretto in modo che le rimozioni garantite non lo consumino. Dopo la correzione,
\str{veto-cover} con tutela attiva raggiunge la medesima riduzione della variante senza tutela
(\SI{17.1}{\percent} su Bridges a soglia 9), con output identico alla configurazione di partenza
su \num{8} configurazioni su \num{8}.

\begin{remark}
Le campagne condotte prima della correzione non sono inficiate: la strategia \str{veto-cover} era
sempre stata impiegata \emph{da sola}, circostanza nella quale l'insieme corrente coincide con
quello in ingresso e i difetti non possono manifestarsi. Il difetto era latente, non realizzato.
\end{remark}

% =======================================================================================
\subsection{D6 --- Indice di rischio anti-correlato}\label{sec:d6}

La politica \str{budget:F} sceglie quali regole rimuovere ordinandole per un indice di rischio
definito come somma dei valori mancanti sulle colonne del lato sinistro. L'indice non era mai
stato validato. La validazione, condotta confrontando l'indice con il numero di celle che ciascuna
regola può effettivamente servire, ha dato una correlazione di rango $\rho = -0.237$ su
\num{8961} regole: l'indice non è semplicemente debole, è \emph{anti-correlato}.

La causa è stata individuata con precisione: l'indice somma i valori mancanti sulle colonne del
lato \emph{sinistro}, mentre una regola può servire soltanto celle della propria colonna di
\emph{destra}. Ne segue che le regole la cui colonna di destra non presenta mai valori mancanti
--- e che quindi non possono servire alcuna cella --- ricevono un indice basso e vengono rimosse
per prime. L'ordinamento promuoveva a bersaglio esattamente le regole innocue.

La correzione adotta un indice calcolato sulle celle nulle reali. A parità di budget, il nuovo
ordinamento rimuove $2.5$ volte più regole che non servono alcuna cella e distrugge $2.6$ volte
meno celle servibili. Come riportato nella \autoref{sec:negative-results}, il miglioramento
strutturale \emph{non} si traduce in un miglioramento dell'esito, e la circostanza ha portato a
concludere che sotto vincolo di budget conti la dimensione del budget e non l'ordinamento.

% =======================================================================================
\section{Che cosa insegnano questi difetti}\label{sec:system-behaviour}

Alcune caratteristiche del sistema non costituiscono difetti ma sono rilevanti per
l'interpretazione dei risultati, e sono raccolte qui insieme alle conseguenze generali del
catalogo.

\paragraph{Codice non raggiungibile.} Il disassemblaggio individua procedure che il codice
applicativo non invoca mai: una procedura di ricerca locale affetta da un errore di segno, due
procedure di riordino delle tuple basate su euristiche, un comparatore di soglie affetto da
troncamento nella conversione a intero, e un meccanismo di gestione di regole privilegiate di
fatto inerte. La circostanza è stata verificata empiricamente oltre che staticamente.

\begin{observation}
Il comparatore di soglie merita una nota, perché il suo troncamento lo rende non transitivo e
avrebbe potuto spiegare l'irregolarità dell'ordine di partenza osservata nella
\autoref{sec:negative-results}. La verifica sui sorgenti decompilati mostra però che nessuna
classe lo istanzia: è codice morto, e non può essere la causa di alcun comportamento osservato.
L'ipotesi è stata scartata su questa base.
\end{observation}

\paragraph{L'argomento di selezione dell'euristica è inerte.} Il sesto argomento della riga di
comando, che dovrebbe selezionare l'euristica di riempimento, non è letto. La circostanza è
stata verificata eseguendo il sistema con valori diversi dell'argomento e confrontando gli
output, che risultano identici. Tutte le osservazioni di questo lavoro valgono per la sola
procedura effettivamente eseguita.

\paragraph{Che cosa avrebbe potuto rilevare prima questi difetti.} Nessuno dei sei difetti è stato
individuato da un test automatico del materiale d'archivio, e per ciascuno esiste una verifica
elementare che lo avrebbe fatto prima. D1 e D4 sono difetti di codice: un test di copertura dei
rami avrebbe mostrato che il ramo categorico non è mai eseguito, e un confronto fra gli eseguibili
distribuiti avrebbe mostrato le due istruzioni di diagnostica. D2 e D3 sono difetti di dati: un
controllo di schema --- il numero di campi per riga e la corrispondenza fra intestazione e
riferimenti --- li avrebbe segnalati al momento della distribuzione, senza bisogno di eseguire
nulla. D5 e D6 sono difetti propri del codice sviluppato in questo lavoro, e in entrambi i casi il
controllo che li ha rivelati è stato un confronto fra ciò che la procedura \emph{dichiara} di fare
e ciò che \emph{fa}: nel primo caso fra riduzione annunciata e riduzione effettiva, nel secondo fra
l'indice di rischio e la grandezza che dovrebbe approssimare. La raccomandazione che se ne trae è
di introdurre questi quattro controlli --- copertura dei rami, confronto fra artefatti, validazione
dello schema, coerenza fra dichiarazione ed effetto --- prima di misurare qualunque cosa su un
sistema opaco.

\paragraph{Che cosa il catalogo insegna sul progetto di algoritmi a RFD.} I sei difetti, letti
insieme, indicano tre proprietà ricorrenti. La prima è la distanza fra documentazione e
comportamento: un sistema può esporre un parametro che non legge, eseguire una sola delle
procedure che dichiara di offrire, e contenere rami che il codice non può raggiungere; per questo
l'analisi del \autoref{chap:renuver} è stata condotta sul bytecode e non sui documenti. La seconda
è la concentrazione del costo: la quasi totalità del tempo è in una procedura sola
(\autoref{sec:cost}), e ciò rende la riduzione dell'insieme di regole l'unica leva temporale di
rilievo, ma anche la più delicata, perché quella stessa procedura è quella che il pruning può
indebolire. La terza riguarda chi misura: due dei sei difetti appartengono agli strumenti
sviluppati in questo lavoro, e la loro individuazione ha richiesto di confrontare in modo
sistematico ciò che una strategia \emph{dichiara} di fare con ciò che \emph{fa}
(\autoref{sec:defects-ours}). Ne segue una raccomandazione operativa per chi impiegasse il sistema
su altri dati: verificare l'effettiva equivalenza degli output prima di attribuire un effetto a
una strategia di potatura, e considerare le configurazioni in cui il file prodotto coincide con
quello di partenza come prove nulle anziché come conferme.
